\section{Introduction}

Whether a community has access to public transit shapes nearly every aspect of daily life: where residents can work, how long commutes take, whether elderly or disabled residents can access healthcare, and what policy priorities communities bring to their legislative representatives. Transit-rich urban communities and transit-poor suburban or rural communities have fundamentally different infrastructure needs, even when they share a zip code or county.

The redistricting criterion to preserve communities of interest implicitly recognises infrastructure access as a community identity marker. Courts have upheld the use of transit-corridor-based community arguments in redistricting cases, particularly in New York and California where transit agencies themselves have filed amicus briefs arguing that transit corridors define community boundaries.

The \textsc{bisect} redistricting pipeline's edge-weight framework provides a natural mechanism for transit accessibility: tracts sharing transit infrastructure receive heavier edges (encouraging METIS to keep them together), while tracts at transit transition boundaries receive lighter edges (inviting cuts at those boundaries). This paper defines the transit accessibility similarity metric, validates it empirically on NC/WI/TX, and demonstrates its legal groundability.

Transit weights (M.7) are the seventh of eight community character signals in the M-track. They complement economic character (M.1), land use (M.2), housing character (M.3), commuting sheds (M.4), topographic separation (M.5), and administrative zones (M.6), and are synthesised in the composite Community Character Index (M.8).
